\documentclass[
    lang=cn,
    scheme=chinese,
    mode=simple,
    device=normal,
    color=blue,
    math=cm,
    thmcnt=section,
    12pt,
    nowatermark
]{elegantbook}

% 项目级水印开关：true 显示，false 关闭。

\newif\ifmathanalysiswatermark
\mathanalysiswatermarkfalse

% =============================================================================
% Project / personal preamble for the integrable-systems notes
% General document-class settings live in elegantbook.cls.
% This file contains only project-specific or frequently adjusted preferences.
% =============================================================================

% ================= Project structure =================
% Kept here rather than in the class because not every ElegantBook document
% needs to be split into subfiles.
\usepackage{subfiles}

% ================= Watermark appearance =================
% Whether the watermark is shown is controlled by the class option:
%   watermark   -> show
%   nowatermark -> hide (default)
% The personal watermark content stays here.
\SetWatermarkText{欧拉的小迷弟}
\SetWatermarkScale{0.4}
\SetWatermarkLightness{0.9}

% ================= Header and footer =================
% Personal layout: current section on the left, nickname on the right,
% page number centered in the footer.
\fancyhf{}
\fancyhead[L]{\rightmark}
\fancyhead[R]{欧拉的小迷弟}
\fancyfoot[C]{\thepage}
\renewcommand{\headrule}{\hrule width\headwidth}
\renewcommand{\headrulewidth}{0.4pt}
\pagestyle{fancy}

% Keep the same personal header/footer on chapter-opening pages.
\fancypagestyle{plain}{%
    \fancyhf{}
    \fancyhead[L]{\rightmark}
    \fancyhead[R]{欧拉的小迷弟}
    \fancyfoot[C]{\thepage}
    \renewcommand{\headrule}{\hrule width\headwidth}
    \renewcommand{\headrulewidth}{0.4pt}%
}

% ================= Obsidian-style supplementary callout =================
% ElegantBook's native remark environment is intentionally left untouched.
% Use callout for supplementary material that should be visually secondary
% during a presentation or review.
\newtcolorbox{callout}[1][]{
    enhanced,
    breakable,
    colback=black!3,
    colframe=black!35,
    boxrule=0pt,
    borderline west={1.4pt}{0pt}{black!45},
    arc=0pt,
    outer arc=0pt,
    left=8pt,
    right=8pt,
    top=7pt,
    bottom=7pt,
    before skip=10pt,
    after skip=10pt,
    title={#1},
    colbacktitle=black!7,
    coltitle=black,
    fonttitle=\bfseries\small,
    fontupper=\small,
    titlerule=0pt
}

% ================= Importance markers =================
\newcommand{\marktriangle}{\ensuremath{\blacktriangle}}
\newcommand{\markstar}{\ensuremath{\bigstar}}
\newcommand{\markstartriangle}{\ensuremath{\bigstar\blacktriangle}}
\newcommand{\marktwostar}{\ensuremath{\bigstar\bigstar}}

% ================= Internal section helper =================
% Preserve unnumbered (A)(B)(C)...-style headings while adding them to the TOC.
\newcommand{\methodsection}[1]{%
    \subsubsection*{#1}%
    \addcontentsline{toc}{subsubsection}{#1}%
}

% ================= Space after examples =================
\newcommand{\exampleblank}{\par\vspace{10\baselineskip}}
\newcommand{\myspace}[1]{\par\vspace{#1\baselineskip}}

% ================= Custom mathematical operators =================
\DeclareMathOperator{\Tr}{Tr}
\DeclareMathOperator{\tr}{tr}
\DeclareMathOperator{\rank}{rank}
\DeclareMathOperator{\diag}{diag}
\DeclareMathOperator{\Ree}{Re}
\DeclareMathOperator{\Imm}{Im}
\DeclareMathOperator{\Ker}{Ker}


\begin{document}

\setcounter{chapter}{4}

\chapter{RH 方法求解零边界 NLS 方程}

\section*{先把整条主线理清楚}

\begin{callout}[前置概念：RH 问题（原书定义 3.12，3.8 节）]
    设 $\Sigma$ 为复平面 $\mathbb{C}$ 内的有向路径，$v(z):\Sigma\to GL(n,\mathbb{C})$ 是 $\Sigma^0$ 上光滑的映射。
    $(\Sigma,v)$ 决定一个 RH 问题：求 $n\times n$ 矩阵 $M(z)$，使它满足

    \begin{itemize}
        \item $M(z)$ 在 $\mathbb{C}\setminus\Sigma$ 上解析（原书 (3.8.1)）；
        \item $M_+(z)=M_-(z)v(z)$，$z\in\Sigma$（原书 (3.8.2)）；
        \item $M(z)\to I$，$z\to\infty$（原书 (3.8.3)）。
    \end{itemize}

    其中 $M_\pm$ 表示在正负域内 $z'$ 趋于 $\Sigma$ 上相应的 $z$ 点时的极限。$\Sigma$ 称为跳跃曲线，$v(z)$ 称为跳跃矩阵。
    根据 Beals-Coifman 定理 (CPAM, 37(1984), 39-90), 如上 RH 问题的解为

    \[
        M(z)
        =
        I+\dfrac{1}{2\pi i}\int\limits_\Sigma\dfrac{\mu(\xi)w(\xi)}{\xi-z}\,\mathrm d\xi,
    \]
    其中 $\mu$ 是奇异积分方程 $(I-C_w)\mu=I$ 的解，$w=w_++w_-$ 由跳跃矩阵的分解 $v=(b_-)^{-1}b_+$ 给出（$w_+=b_+-I$、$w_-=I-b_-$）。算子 $C_w$ 定义为

    \[
        C_wf=C_+(fw_-)+C_-(fw_+),
    \]
    而 $C_\pm$ 是 Cauchy 投影

    \[
        (C_\pm f)(z)=\lim\limits_{\substack{z'\to z\\ z'\in\pm\Sigma}}\dfrac{1}{2\pi i}\int\limits_\Sigma\dfrac{f(\xi)}{\xi-z'}\,\mathrm d\xi.
    \]
    以上即原书 (3.8.6)---(3.8.10)。
\end{callout}

\section{NLS 方程、Lax 对与大谱参数重构}

\subsection{NLS 初值问题}

考虑零边界聚焦 NLS 方程

\begin{equation}
    iq_t(x,t)+q_{xx}(x,t)+2q(x,t)\lvert q(x,t)\rvert^2=0.
    \tag{5.1.1}
    \label{eq:5.1.1}
\end{equation}

初值为

\begin{equation}
    q(x,0)=q_0(x)\in\mathcal S(\mathbb{R}).
    \tag{5.1.2}
    \label{eq:5.1.2}
\end{equation}


\subsection{Lax 对}

NLS 方程具有矩阵形式的 Lax 对

\begin{equation}
    \psi_x+iz\sigma_3\psi=P\psi,
    \tag{5.1.3}
    \label{eq:5.1.3}
\end{equation}

\begin{equation}
    \psi_t+2iz^2\sigma_3\psi=Q\psi,
    \tag{5.1.4}
    \label{eq:5.1.4}
\end{equation}

其中

\[
    \sigma_3=
    \begin{pmatrix}
        1 & 0  \\
        0 & -1
    \end{pmatrix},
    \qquad
    P=
    \begin{pmatrix}
        0       & q \\
        -\bar q & 0
    \end{pmatrix},
\]

\[
    Q=
    \begin{pmatrix}
        i\lvert q\rvert^2 & iq_x               \\
        i\bar q_x         & -i\lvert q\rvert^2
    \end{pmatrix}
    +2zP.
\]

\begin{callout}[前置概念：Lax 对、谱参数和零曲率条件]
    把 \eqref{eq:5.1.3}---\eqref{eq:5.1.4} 写成

    \[
        \psi_x=U\psi,\qquad \psi_t=V\psi,
    \]
    其中 $U=P-iz\sigma_3$、$V=Q-2iz^2\sigma_3$。因为同一个 $\psi$ 同时满足两式，所以必须有

    \[
        \psi_{xt}=\psi_{tx}.
    \]
    展开后得到

    \[
        U_t-V_x+[U,V]=0.
    \]
    这就是零曲率条件。把这里的 $U,V$ 代入，非对角元恰好恢复 NLS \eqref{eq:5.1.1}。因此 Lax 对把非线性方程编码成一对关于 $\psi$ 的线性系统。
    $z$ 称为谱参数。固定 $t,z$ 后，\eqref{eq:5.1.3} 是关于 $x$ 的一阶线性矩阵微分方程。
\end{callout}

\begin{callout}[计算：零曲率条件怎样恢复 NLS？]
    写成

    \[
        \psi_x=U\psi,\qquad
        \psi_t=V\psi,
    \]
    其中

    \[
        U=
        P-iz\sigma_3
        =
        \begin{pmatrix}
            -iz     & q  \\
            -\bar q & iz
        \end{pmatrix},
    \]

    \[
        V=
        Q-2iz^2\sigma_3
        =
        \begin{pmatrix}
            i|q|^2-2iz^2       & iq_x+2zq      \\
            i\bar q_x-2z\bar q & -i|q|^2+2iz^2
        \end{pmatrix}.
    \]
    计算可得

    \[
        U_t-V_x+[U,V]
        =
        \begin{pmatrix}
            0                                    &
            q_t-iq_{xx}-2i|q|^2q                   \\
            -\bar q_t-i\bar q_{xx}-2i|q|^2\bar q &
            0
        \end{pmatrix}.
    \]
    令其为零，右上角给出

    \[
        q_t-iq_{xx}-2i|q|^2q=0.
    \]
    两边乘 $i$：

    \[
        iq_t+q_{xx}+2q|q|^2=0,
    \]
    正好就是 NLS 方程 \eqref{eq:5.1.1}；左下角是它的共轭方程。
\end{callout}

\subsection{Jost 解与变换}

当 $\lvert x\rvert\to\infty$ 时，$q\to0$，所以 $P\to0$，同时 $Q$ 的非自由部分也趋于零。Lax 对渐近为

\[
    \psi_x+iz\sigma_3\psi=0,\qquad
    \psi_t+2iz^2\sigma_3\psi=0.
\]

自由基本矩阵解为

\[
    e^{-i\theta(z)\sigma_3},
    \qquad
    \theta(z)=zx+2z^2t.
\]

Jost 解是在空间无穷远处与这个自由解匹配的特殊基本矩阵解。为了剥离已知的自由振荡，作变换

\begin{equation}
    \mu(x,t,z)=\psi(x,t,z)e^{i\theta(z)\sigma_3}.
    \tag{5.1.5}
    \label{eq:5.1.5}
\end{equation}

因此

\begin{equation}
    \mu(x,t,z)\sim I,\qquad \lvert x\rvert\to\infty.
    \tag{5.1.6}
    \label{eq:5.1.6}
\end{equation}

\begin{callout}[计算：这个自由基本解是怎样求出来的？]
    先固定 $t,z$ 解 $x$-方程：

    \[
        \psi_x=-iz\sigma_3\psi.
    \]
    由于 $\sigma_3$ 是常矩阵，

    \[
        \psi(x,t,z)
        =
        e^{-izx\sigma_3}C(t,z),
    \]
    其中 $C(t,z)$ 与 $x$ 无关。再代入 $t$-方程

    \[
        \psi_t+2iz^2\sigma_3\psi=0.
    \]
    因为

    \[
        \psi_t=e^{-izx\sigma_3}C_t,
    \]
    所以

    \[
        C_t+2iz^2\sigma_3C=0,
    \]
    从而

    \[
        C(t,z)
        =
        e^{-2iz^2t\sigma_3}C_0(z).
    \]
    于是

    \[
        \psi(x,t,z)
        =
        e^{-i(zx+2z^2t)\sigma_3}C_0(z)
        =
        e^{-i\theta(z)\sigma_3}C_0(z).
    \]
    选择 Jost 归一化时取相应的 $C_0(z)=I$，所以自由归一化解为

    \[
        e^{-i\theta(z)\sigma_3}
        =
        \begin{pmatrix}
            e^{-i\theta(z)} & 0              \\
            0               & e^{i\theta(z)}
        \end{pmatrix}.
    \]

\end{callout}

\begin{callout}[概念：Jost 解为什么有左右两个？]
    固定 $t,z$ 后，Lax 对的 $x$ 部分是一个 $2\times2$ 一阶线性系统。
    一个基本矩阵解本身包含两个线性无关的列解，而且基本矩阵解并不唯一：若 $\Psi$ 是基本矩阵解，则 $\Psi C$（$C$ 为可逆常矩阵）仍是基本矩阵解。
    整条实轴有两个自然的归一化端点，所以分别选取

    \[
        \psi_1e^{i\theta\sigma_3}\to I,\qquad x\to-\infty,
    \]

    \[
        \psi_2e^{i\theta\sigma_3}\to I,\qquad x\to+\infty.
    \]
\end{callout}

\subsection{等价 Lax 对与全微分形式}

由 \eqref{eq:5.1.5} 得到与原 Lax 对等价的系统

\begin{equation}
    \mu_x+iz[\sigma_3,\mu]=P\mu,
    \tag{5.1.7}
    \label{eq:5.1.7}
\end{equation}

\begin{equation}
    \mu_t+2iz^2[\sigma_3,\mu]=Q\mu.
    \tag{5.1.8}
    \label{eq:5.1.8}
\end{equation}

进一步写成全微分形式

\begin{equation}
    \mathrm d\!\left(e^{i\theta(z)\widehat{\sigma}_3}\mu\right)
    =
    e^{i\theta(z)\widehat{\sigma}_3}
    \left[(P\,\mathrm dx+Q\,\mathrm dt)\mu\right].
    \tag{5.1.9}
    \label{eq:5.1.9}
\end{equation}

\begin{callout}[计算：从 $\psi$ 的 Lax 对得到 \eqref{eq:5.1.7}---\eqref{eq:5.1.9}]
    由 \eqref{eq:5.1.5}

    \[
        \psi=\mu e^{-i\theta\sigma_3}.
    \]
    对 $x$ 求导，利用 $\theta_x=z$：

    \[
        \psi_x
        =
        \left(
        \mu_x-iz\mu\sigma_3
        \right)e^{-i\theta\sigma_3}.
    \]
    代入 \eqref{eq:5.1.3}：

    \[
        \left(
        \mu_x-iz\mu\sigma_3
        \right)e^{-i\theta\sigma_3}
        +
        iz\sigma_3\mu e^{-i\theta\sigma_3}
        =
        P\mu e^{-i\theta\sigma_3}.
    \]
    右乘 $e^{i\theta\sigma_3}$，得到

    \[
        \mu_x+iz(\sigma_3\mu-\mu\sigma_3)=P\mu,
    \]
    即 \eqref{eq:5.1.7}。同理利用 $\theta_t=2z^2$ 得到 \eqref{eq:5.1.8}。
    再令

    \[
        F=e^{i\theta\widehat{\sigma}_3}\mu.
    \]
    则

    \[
        F_x
        =
        e^{i\theta\widehat{\sigma}_3}
        \left(
        \mu_x+iz[\sigma_3,\mu]
        \right)
        =
        e^{i\theta\widehat{\sigma}_3}P\mu,
    \]

    \[
        F_t
        =
        e^{i\theta\widehat{\sigma}_3}
        \left(
        \mu_t+2iz^2[\sigma_3,\mu]
        \right)
        =
        e^{i\theta\widehat{\sigma}_3}Q\mu.
    \]
    因此

    \[
        \mathrm dF
        =
        F_x\,\mathrm dx+F_t\,\mathrm dt
        =
        e^{i\theta\widehat{\sigma}_3}
        \left[(P\,\mathrm dx+Q\,\mathrm dt)\mu\right],
    \]
    这就是 \eqref{eq:5.1.9}。
\end{callout}

\begin{callout}[概念：$\widehat{\sigma}_3$ 是什么？]
    $\widehat{\sigma}_3$是作用在矩阵空间上的交换子算子：

    \[
        \widehat{\sigma}_3A=[\sigma_3,A]=\sigma_3A-A\sigma_3.
    \]
    其指数满足

    \[
        e^{a\widehat{\sigma}_3}A
        =
        e^{a\sigma_3}Ae^{-a\sigma_3}.
    \]
    若

    \[
        A=\begin{pmatrix}a&b\\c&d\end{pmatrix},
    \]
    则

    \[
        e^{\alpha\widehat{\sigma}_3}A
        =
        \begin{pmatrix}
            a             & e^{2\alpha}b \\
            e^{-2\alpha}c & d
        \end{pmatrix}.
    \]
    因而这个算子真正改变的是非对角元，这正是后面出现 $e^{\pm2iz(x-y)}$ 的原因。
\end{callout}

\subsection{为什么还要证明 \texorpdfstring{$\mu\to I$}{μ → I}（\texorpdfstring{$z\to\infty$}{z → ∞}）？}

构造规范 RH 问题时，希望 RH 未知函数满足

\[
    M(z)\to I,\qquad z\to\infty.
\]

所以除了已经有的空间无穷远归一化 \eqref{eq:5.1.6}，还要证明谱参数无穷远处

\[
    \mu(x,t,z)\to I,\qquad z\to\infty.
\]

将 $\mu$ 在无穷远点作渐近展开

\begin{equation}
    \mu
    =
    \mu^{(0)}
    +\dfrac{\mu^{(1)}}{z}
    +\dfrac{\mu^{(2)}}{z^2}
    +\cdots
    =
    \mu^{(0)}+O(z^{-1}),
    \qquad z\to\infty.
    \tag{5.1.10}
    \label{eq:5.1.10}
\end{equation}

其中 $\mu^{(0)},\mu^{(1)},\ldots$ 与 $z$ 无关。下面把 \eqref{eq:5.1.10} 分别代入 Lax 对的 $x$-部分和 $t$-部分，并按 $z$ 的幂次比较系数。

先看 $x$-部分 \eqref{eq:5.1.7}：

\[
    \mu_x+iz[\sigma_3,\mu]=P\mu.
\]

由 \eqref{eq:5.1.10}，有

\[
    \mu_x
    =
    \mu_x^{(0)}
    +\dfrac{\mu_x^{(1)}}{z}
    +\dfrac{\mu_x^{(2)}}{z^2}
    +\cdots,
\]

\[
    iz[\sigma_3,\mu]
    =
    iz[\sigma_3,\mu^{(0)}]
    +i[\sigma_3,\mu^{(1)}]
    +\dfrac{i[\sigma_3,\mu^{(2)}]}{z}
    +\cdots,
\]

而右端

\[
    P\mu
    =
    P\mu^{(0)}
    +\dfrac{P\mu^{(1)}}{z}
    +\dfrac{P\mu^{(2)}}{z^2}
    +\cdots.
\]

因此整个 $x$-部分展开为

\[
    iz[\sigma_3,\mu^{(0)}]
    +\Bigl(\mu_x^{(0)}+i[\sigma_3,\mu^{(1)}]\Bigr)
    +\dfrac{\mu_x^{(1)}+i[\sigma_3,\mu^{(2)}]}{z}
    +\cdots
    =
    P\mu^{(0)}
    +\dfrac{P\mu^{(1)}}{z}
    +\cdots.
\]

比较 $O(z)$ 项，得到

\begin{equation}
    O(z):\quad[\sigma_3,\mu^{(0)}]=0,
    \tag{5.1.11}
    \label{eq:5.1.11}
\end{equation}

比较 $O(z^0)$ 项，得到

\begin{equation}
    O(z^0):\quad\mu_x^{(0)}+i[\sigma_3,\mu^{(1)}]=P\mu^{(0)}.
    \tag{5.1.12}
    \label{eq:5.1.12}
\end{equation}

再看 $t$-部分 \eqref{eq:5.1.8}：

\[
    \mu_t+2iz^2[\sigma_3,\mu]=Q\mu.
\]

将 $Q$ 分成关于 $z$ 的两部分：

\[
    Q=Q_0+2zP,
    \qquad
    Q_0=
    \begin{pmatrix}
        i\lvert q\rvert^2 & iq_x \\
        i\bar q_x          & -i\lvert q\rvert^2
    \end{pmatrix}.
\]

于是左端展开为

\[
    \mu_t+2iz^2[\sigma_3,\mu]
    =
    2iz^2[\sigma_3,\mu^{(0)}]
    +2iz[\sigma_3,\mu^{(1)}]
    +\Bigl(\mu_t^{(0)}+2i[\sigma_3,\mu^{(2)}]\Bigr)
    +O(z^{-1}),
\]

右端展开为

\[
    \begin{aligned}
        Q\mu
        &=(Q_0+2zP)
        \left(
            \mu^{(0)}+\dfrac{\mu^{(1)}}{z}+\dfrac{\mu^{(2)}}{z^2}+\cdots
        \right)\\
        &=2zP\mu^{(0)}
        +\Bigl(Q_0\mu^{(0)}+2P\mu^{(1)}\Bigr)
        +O(z^{-1}).
    \end{aligned}
\]

比较 $O(z^2)$ 项，得到

\begin{equation}
    O(z^2):\quad[\sigma_3,\mu^{(0)}]=0,
    \tag{5.1.13}
    \label{eq:5.1.13}
\end{equation}

比较 $O(z)$ 项，得到

\[
    2i[\sigma_3,\mu^{(1)}]=2P\mu^{(0)},
\]

即

\begin{equation}
    O(z):\quad i[\sigma_3,\mu^{(1)}]=P\mu^{(0)}.
    \tag{5.1.14}
    \label{eq:5.1.14}
\end{equation}

下面说明为什么 \eqref{eq:5.1.11}、\eqref{eq:5.1.13} 能推出 $\mu^{(0)}$ 为对角矩阵。设

\[
    \mu^{(0)}=
    \begin{pmatrix}
        a & b \\
        c & d
    \end{pmatrix}.
\]

则

\[
    [\sigma_3,\mu^{(0)}]
    =
    \begin{pmatrix}
        1 & 0 \\
        0 & -1
    \end{pmatrix}
    \begin{pmatrix}
        a & b \\
        c & d
    \end{pmatrix}
    -
    \begin{pmatrix}
        a & b \\
        c & d
    \end{pmatrix}
    \begin{pmatrix}
        1 & 0 \\
        0 & -1
    \end{pmatrix}
    =
    \begin{pmatrix}
        0  & 2b \\
        -2c & 0
    \end{pmatrix}.
\]

由 $[\sigma_3,\mu^{(0)}]=0$ 得 $b=c=0$，所以

\[
    \mu^{(0)}=
    \begin{pmatrix}
        a & 0 \\
        0 & d
    \end{pmatrix}
\]

确实是对角矩阵。另一方面，将 \eqref{eq:5.1.14} 代入 \eqref{eq:5.1.12}，或者直接用两式相减，有

\[
    \mu_x^{(0)}
    +\underbrace{i[\sigma_3,\mu^{(1)}]}_{P\mu^{(0)}}
    =P\mu^{(0)},
\]

因此

\[
    \mu_x^{(0)}=0.
\]

这说明 $\mu^{(0)}$ 是与 $x$ 无关的对角矩阵。因此 \eqref{eq:5.1.10} 对 $x,z$ 同时取极限，并注意交换极限顺序，我们有

\[
    \lim\limits_{z\to\infty}\lim\limits_{\lvert x\rvert\to\infty}\mu
    =
    \lim\limits_{\lvert x\rvert\to\infty}\lim\limits_{z\to\infty}
    \left(\mu^{(0)}+O(z^{-1})\right).
\]

利用 \eqref{eq:5.1.6} 和 \eqref{eq:5.1.10}，上述左边极限为 $I$，右边为 $\mu^{(0)}$，因此

\[
    \mu^{(0)}=I,
\]

从而

\[
    \mu\to I,\qquad z\to\infty.
\]

最后将 $\mu^{(0)}=I$ 代回 \eqref{eq:5.1.14}：

\[
    i[\sigma_3,\mu^{(1)}]=P.
\]

若记

\[
    \mu^{(1)}=
    \begin{pmatrix}
        \mu_{11}^{(1)} & \mu_{12}^{(1)} \\
        \mu_{21}^{(1)} & \mu_{22}^{(1)}
    \end{pmatrix},
\]

则

\[
    i[\sigma_3,\mu^{(1)}]
    =
    \begin{pmatrix}
        0 & 2i\mu_{12}^{(1)} \\
        -2i\mu_{21}^{(1)} & 0
    \end{pmatrix}
    =
    P
    =
    \begin{pmatrix}
        0 & q \\
        -\bar q & 0
    \end{pmatrix}.
\]

比较 $(1,2)$ 元，得到

\begin{equation}
    q(x,t)
    =
    2i(\mu^{(1)})_{12}
    =
    2i\lim\limits_{z\to\infty}(z\mu)_{12}.
    \tag{5.1.15}
    \label{eq:5.1.15}
\end{equation}

其中下标 $12$ 表示该矩阵的第一行第二列元素。

\begin{callout}[式 \eqref{eq:5.1.15} 的作用]
    \eqref{eq:5.1.15} 第一次把 \textbf{NLS 的解 $q$} 与 \textbf{特征函数 $\mu$} 联系起来。接下来第 2、3 节要继续把特征函数与 RH 问题联系起来。这样最终就形成

    \[
        \text{NLS 解 }q
        \longleftrightarrow
        \text{特征函数}
        \longleftrightarrow
        \text{RH 问题的解}.
    \]
    因而以后不再直接求 NLS，而是求 RH 问题，再用 \eqref{eq:5.1.15} 或其 RH 版本恢复 $q$。
\end{callout}

\section{特征函数的解析性与对称性}

由于 \eqref{eq:5.1.9} 是全微分形式，在相应区域中积分只与端点有关。取固定 $t$ 的两条水平路径

\[
    (-\infty,t)\to(x,t),
    \qquad
    (+\infty,t)\to(x,t),
\]

得到两个特征函数

\begin{equation}
    \mu_1(x,t,z)
    =
    I+
    \int_{-\infty}^{x}
    e^{-iz(x-y)\widehat{\sigma}_3}
    P(y,t)\mu_1(y,t,z)\,\mathrm dy,
    \tag{5.2.1}
    \label{eq:5.2.1}
\end{equation}

\begin{equation}
    \mu_2(x,t,z)
    =
    I-
    \int_x^{\infty}
    e^{-iz(x-y)\widehat{\sigma}_3}
    P(y,t)\mu_2(y,t,z)\,\mathrm dy.
    \tag{5.2.2}
    \label{eq:5.2.2}
\end{equation}

它们满足各自的空间归一化

\[
    \mu_1\to I\quad(x\to-\infty),
    \qquad
    \mu_2\to I\quad(x\to+\infty),
\]

并且由第 5.1 节的大 $z$ 分析，

\[
    \mu_1,\mu_2\to I,\qquad z\to\infty.
\]

\begin{callout}[概念：为什么 \eqref{eq:5.2.1}---\eqref{eq:5.2.2} 是 Volterra 积分方程？（原书 5.2.1 节）]
    Volterra 型积分方程的积分区间端点包含当前点 $x$，例如 $\displaystyle\int_{-\infty}^x$ 或 $\displaystyle\int_x^\infty$。
    这种结构使它适合逐次迭代，后面存在性与唯一性正是通过 Neumann 级数和 Bellman/Grönwall 型估计得到的。
\end{callout}

令

\[
    \psi_1=\mu_1e^{-i\theta(z)\sigma_3},
    \qquad
    \psi_2=\mu_2e^{-i\theta(z)\sigma_3}.
\]

二者都是原 Lax 对 \eqref{eq:5.1.3}---\eqref{eq:5.1.4} 的基本矩阵解，所以存在与 $x,t$ 无关的可逆矩阵 $S(z)$，使$\psi_1=\psi_2S$，即

\begin{equation}
    \mu_1(x,t,z)
    =
    \mu_2(x,t,z)e^{-i\theta(z)\widehat{\sigma}_3}S(z).
    \tag{5.2.3}
    \label{eq:5.2.3}
\end{equation}

其中

\[
    S(z)=
    \begin{pmatrix}
        s_{11}(z) & s_{12}(z) \\
        s_{21}(z) & s_{22}(z)
    \end{pmatrix}
\]

称为谱矩阵或散射矩阵。

\begin{callout}[为什么 $S$ 与 $x,t$ 无关？]
    原始 Jost 解满足 $\psi_1=\psi_2S$，即 $S=\psi_2^{-1}\psi_1$。对 $x$ 求导：

    \[
        S_x
        =
        -\psi_2^{-1}U\psi_1+\psi_2^{-1}U\psi_1
        =0.
    \]
    对 $t$ 同理有 $S_t=0$。因此在本书采用的规范下，$S=S(z)$；时空依赖进入相位 $\theta$、跳跃矩阵和离散核向量中。
\end{callout}

\subsection{解析性}

记 $\mu_1$ 的两列为

\begin{equation}
    \mu_1
    =
    \begin{pmatrix}
        \mu_1^{(11)} & \mu_1^{(12)} \\
        \mu_1^{(21)} & \mu_1^{(22)}
    \end{pmatrix}
    =
    (\mu_1^+,\mu_1^-).
    \tag{5.2.4}
    \label{eq:5.2.4}
\end{equation}

由 \eqref{eq:5.2.1} 分列得到

\begin{equation}
    \mu_1^+(x,t,z)
    =
    \begin{pmatrix}1\\0\end{pmatrix}
    +
    \int_{-\infty}^{x}
    \operatorname{diag}\!\left(1,e^{2iz(x-y)}\right)
    P(y,t)\mu_1^+(y,t,z)\,\mathrm dy,
    \tag{5.2.5}
    \label{eq:5.2.5}
\end{equation}

\begin{equation}
    \mu_1^-(x,t,z)
    =
    \begin{pmatrix}0\\1\end{pmatrix}
    +
    \int_{-\infty}^{x}
    \operatorname{diag}\!\left(e^{-2iz(x-y)},1\right)
    P(y,t)\mu_1^-(y,t,z)\,\mathrm dy.
    \tag{5.2.6}
    \label{eq:5.2.6}
\end{equation}

对 \eqref{eq:5.2.5}---\eqref{eq:5.2.6}，积分区域满足 $y<x$，即 $x-y>0$。令 $z=\operatorname{Re}z+i\operatorname{Im}z$，则

\[
    e^{2iz(x-y)}
    =
    e^{2i(x-y)\operatorname{Re}z}
    e^{-2(x-y)\operatorname{Im}z},
\]

\[
    e^{-2iz(x-y)}
    =
    e^{-2i(x-y)\operatorname{Re}z}
    e^{2(x-y)\operatorname{Im}z}.
\]

因此：

\begin{itemize}
    \item 当 $\operatorname{Im}z>0$ 时，$e^{2iz(x-y)}$ 衰减，所以 $\mu_1^+$ 在 $\mathbb{C}^+$ 解析；
    \item 当 $\operatorname{Im}z<0$ 时，$e^{-2iz(x-y)}$ 衰减，所以 $\mu_1^-$ 在 $\mathbb{C}^-$ 解析。
\end{itemize}

\begin{callout}[存在性、唯一性与解析性的严格证明思路（原书 5.2.1 节）]
    书中对 $\mu_1^+$ 用 Neumann 级数作严格证明。令

    \begin{equation}
        w^{(0)}=\begin{pmatrix}1\\0\end{pmatrix},
        \qquad
        w^{(n+1)}(x,t,z)
        =
        \int_{-\infty}^{x}
        F(x,y,z)w^{(n)}(y,t,z)\,\mathrm dy,
        \tag{5.2.7}
        \label{eq:5.2.7}
    \end{equation}
    其中

    \begin{equation}
        F(x,y,z)
        =
        \operatorname{diag}\!\left(1,e^{2iz(x-y)}\right)P(y,t).
        \tag{5.2.8}
        \label{eq:5.2.8}
    \end{equation}
    构造 Neumann 级数

    \begin{equation}
        \sum\limits_{n=0}^{\infty}w^{(n)}(x,t,z).
        \tag{5.2.9}
        \label{eq:5.2.9}
    \end{equation}
    取向量 $L^1$ 范数，有

    \begin{equation}
        \lVert w^{(n+1)}(x,t,z)\rVert
        \le
        \int_{-\infty}^{x}
        \lVert F\rVert
        \lVert w^{(n)}(y,t,z)\rVert\,\mathrm dy,
        \tag{5.2.10}
        \label{eq:5.2.10}
    \end{equation}
    且在相应半平面

    \begin{equation}
        \lVert F\rVert\le4\lvert q\rvert.
        \tag{5.2.11}
        \label{eq:5.2.11}
    \end{equation}
    因此

    \begin{equation}
        \lVert w^{(n+1)}(x,t,z)\rVert
        \le
        \int_{-\infty}^{x}
        4\lvert q\rvert
        \lVert w^{(n)}(y,t,z)\rVert\,\mathrm dy.
        \tag{5.2.12}
        \label{eq:5.2.12}
    \end{equation}
    记

    \begin{equation}
        \rho(x,t)
        =
        \int_{-\infty}^{x}4\lvert q(y,t)\rvert\,\mathrm dy,
        \tag{5.2.13}
        \label{eq:5.2.13}
    \end{equation}
    则

    \begin{equation}
        \rho_x=4\lvert q\rvert,
        \qquad
        \lVert w^{(1)}\rVert\le\rho(x,t).
        \tag{5.2.14}
        \label{eq:5.2.14}
    \end{equation}
    归纳得到

    \begin{equation}
        \lVert w^{(n)}(x,t,z)\rVert
        \le
        \dfrac{\rho^n(x,t)}{n!}.
        \tag{5.2.15}
        \label{eq:5.2.15}
    \end{equation}
    对 $x\le a$，令 $\sigma=\displaystyle\int_{-\infty}^{a}4\lvert q\rvert\,\mathrm dy$，则

    \begin{equation}
        \lVert w^{(n)}(x,t,z)\rVert
        \le
        \dfrac{\sigma^n}{n!}.
        \tag{5.2.16}
        \label{eq:5.2.16}
    \end{equation}
    因而级数绝对一致收敛，定义

    \begin{equation}
        w(x,t,z)
        =
        \sum\limits_{n=0}^{\infty}w^{(n)}(x,t,z),
        \tag{5.2.17}
        \label{eq:5.2.17}
    \end{equation}
    并满足

    \begin{equation}
        w(x,t,z)
        =
        w^{(0)}
        +
        \int_{-\infty}^{x}
        F(x,y,z)w(y,t,z)\,\mathrm dy.
        \tag{5.2.18}
        \label{eq:5.2.18}
    \end{equation}
    每个迭代项对 $z$ 解析，一致收敛保证极限仍解析；唯一性则对两个解之差使用 Bellman/Grönwall 型不等式。
    组会上不需要展开证明，说明``Volterra 迭代 + 阶乘估计 + 一致收敛 + 唯一性估计''即可。
\end{callout}

同理，

\begin{equation}
    \mu_2
    =
    \begin{pmatrix}
        \mu_2^{(11)} & \mu_2^{(12)} \\
        \mu_2^{(21)} & \mu_2^{(22)}
    \end{pmatrix}
    =
    (\mu_2^-,\mu_2^+),
    \tag{5.2.19}
    \label{eq:5.2.19}
\end{equation}

其中 $\mu_2^-$ 在 $\mathbb{C}^-$ 解析，$\mu_2^+$ 在 $\mathbb{C}^+$ 解析。
接下来要得到逆矩阵各行的解析性，因此先证明 $\det\mu_j=1$。注意 $\psi_1,\psi_2$ 均满足 Lax 对，并且

具体地，

\[
    \operatorname{tr}(P-iz\sigma_3)
    =
    \operatorname{tr}
    \begin{pmatrix}
        -iz & q \\
        -\bar q & iz
    \end{pmatrix}
    =
    (-iz)+(iz)
    =0,
\]

而

\[
    \operatorname{tr}(Q-2iz^2\sigma_3)
    =
    \operatorname{tr}
    \begin{pmatrix}
        i\lvert q\rvert^2-2iz^2 & iq_x+2zq \\
        i\bar q_x-2z\bar q & -i\lvert q\rvert^2+2iz^2
    \end{pmatrix}
    =0.
\]

由 Abel 公式，

\begin{equation}
    (\det\psi_j)_x
    =
    (\det\psi_j)_t
    =
    0.
    \tag{5.2.20}
    \label{eq:5.2.20}
\end{equation}

又因为 $\mu_j=\psi_je^{i\theta\sigma_3}$，并且

\[
    e^{i\theta\sigma_3}
    =
    \begin{pmatrix}
        e^{i\theta} & 0 \\
        0 & e^{-i\theta}
    \end{pmatrix},
\]

所以

\[
    \det e^{i\theta\sigma_3}
    =
    e^{i\theta}e^{-i\theta}
    =1.
\]

因此 $\det\mu_j=\det\psi_j$。结合空间无穷远归一化，

\begin{equation}
    \det\mu_j=1,\qquad j=1,2.
    \tag{5.2.21}
    \label{eq:5.2.21}
\end{equation}

从 \eqref{eq:5.2.3} 再取行列式，得到

\[
    \det S(z)=1.
\]

\begin{callout}[前置工具：Abel 公式（原书第 2 章定理 2.12）]
    对矩阵微分方程

    \[
        Y_x=A(x)Y,
    \]
    有

    \[
        (\det Y)_x=\operatorname{tr}(A)\det Y,
    \]
    从而

    \[
        \det Y(x)
        =
        \det Y(x_0)
        \exp\!\left(\int_{x_0}^{x}\operatorname{tr}A(s)\,\mathrm ds\right).
    \]
    特别地，若 $\operatorname{tr}A=0$，则 $\det Y$ 为常数。本节只需要这个零迹特例。
\end{callout}

由于 $\det\mu_j=1$，$\mu_1,\mu_2$ 可逆，并且对 $2\times2$ 矩阵逆矩阵可直接写成伴随矩阵。具体地，如果

\[
    \mu=
    \begin{pmatrix}
        a & b \\
        c & d
    \end{pmatrix},
    \qquad
    \det\mu=ad-bc=1,
\]

那么一般的 $2\times2$ 逆矩阵公式给出

\[
    \mu^{-1}
    =
    \dfrac{1}{ad-bc}
    \begin{pmatrix}
        d  & -b \\
        -c & a
    \end{pmatrix}
    =
    \begin{pmatrix}
        d  & -b \\
        -c & a
    \end{pmatrix}.
\]

这就是为什么 $\det\mu_j=1$ 后，逆矩阵的元素可以直接由原矩阵的元素写出。特别地，如果 $\mu_1$ 的第一列在 $\mathbb{C}^+$ 解析、第二列在 $\mathbb{C}^-$ 解析，那么

\[
    \mu_1^{-1}
    =
    \begin{pmatrix}
        \mu_1^{(22)}  & -\mu_1^{(12)} \\
        -\mu_1^{(21)} & \mu_1^{(11)}
    \end{pmatrix}
\]

的第一行全部来自 $\mu_1$ 的第二列，所以在 $\mathbb{C}^-$ 解析；第二行全部来自第一列，所以在 $\mathbb{C}^+$ 解析。于是

\begin{equation}
    \mu_1^{-1}
    =
    \begin{pmatrix}
        \mu_1^{(22)}  & -\mu_1^{(12)} \\
        -\mu_1^{(21)} & \mu_1^{(11)}
    \end{pmatrix}
    =
    \begin{pmatrix}
        \widehat\mu_1^- \\
        \widehat\mu_1^+
    \end{pmatrix},
    \tag{5.2.22}
    \label{eq:5.2.22}
\end{equation}

\begin{equation}
    \mu_2^{-1}
    =
    \begin{pmatrix}
        \mu_2^{(22)}  & -\mu_2^{(12)} \\
        -\mu_2^{(21)} & \mu_2^{(11)}
    \end{pmatrix}
    =
    \begin{pmatrix}
        \widehat\mu_2^+ \\
        \widehat\mu_2^-
    \end{pmatrix}.
    \tag{5.2.23}
    \label{eq:5.2.23}
\end{equation}

这里帽号 $\widehat\mu_j^\pm$ 用来表示\textbf{行向量}。
由 \eqref{eq:5.2.22}---\eqref{eq:5.2.23} 可见：$\mu_1^{-1}$ 第一、二行分别在 $\mathbb{C}^-$、$\mathbb{C}^+$ 解析；$\mu_2^{-1}$ 第一、二行分别在 $\mathbb{C}^+$、$\mathbb{C}^-$ 解析。
由 \eqref{eq:5.2.3} 得

\[
    \begin{aligned}
        e^{-i\theta(z)\widehat{\sigma}_3}S(z)
        & =
        e^{-i\theta(z)\sigma_3}S(z)e^{i\theta(z)\sigma_3} \\
        & =
        \begin{pmatrix}
            s_{11}(z) & s_{12}(z)e^{-2i\theta(z)} \\
            s_{21}(z)e^{2i\theta(z)} & s_{22}(z)
        \end{pmatrix} \\
        & =
        \mu_2^{-1}\mu_1 \\
        & =
        \begin{pmatrix}
            \widehat\mu_2^+ \\
            \widehat\mu_2^-
        \end{pmatrix}
        (\mu_1^+,\mu_1^-) \\
        & =
        \begin{pmatrix}
            \widehat\mu_2^+\mu_1^+ & \widehat\mu_2^+\mu_1^- \\
            \widehat\mu_2^-\mu_1^+ & \widehat\mu_2^-\mu_1^-
        \end{pmatrix}.
    \end{aligned}
\]

因此对角元的解析性直接读出：

\[
    s_{11}(z)\text{ 在 }\mathbb{C}^+\text{ 解析},
    \qquad
    s_{22}(z)\text{ 在 }\mathbb{C}^-\text{ 解析}.
\]

而 $s_{12},s_{21}$ 一般不能解析延拓到整个上、下半平面，但连续到连续谱所在的实轴 $\mathbb{R}$。

\subsection{对称性}

这一部分主要掌握证明思路和最终结论。原 $x$-部分为

\begin{equation}
    \mu_{j,x}(x,t,z)
    +
    iz[\sigma_3,\mu_j(x,t,z)]
    =
    P\mu_j(x,t,z),
    \qquad j=1,2.
    \tag{5.2.24}
    \label{eq:5.2.24}
\end{equation}

把 $z$ 换成 $\bar z$ 后取 Hermite 共轭，并用 $P^H=-P$，得到

\begin{equation}
    \mu_{j,x}^H(x,t,\bar z)
    +
    iz[\sigma_3,\mu_j^H(x,t,\bar z)]
    =
    -\mu_j^H(x,t,\bar z)P.
    \tag{5.2.25}
    \label{eq:5.2.25}
\end{equation}

另一方面，由 $\mu_j\mu_j^{-1}=I$ 对 $x$ 求导：

\begin{equation}
    (\mu_j^{-1})_x
    =
    -\mu_j^{-1}\mu_{j,x}\mu_j^{-1}.
    \tag{5.2.26}
    \label{eq:5.2.26}
\end{equation}

代入 \eqref{eq:5.2.24}，整理为

\begin{equation}
    (\mu_j^{-1})_x
    +
    iz[\sigma_3,\mu_j^{-1}]
    =
    -\mu_j^{-1}P.
    \tag{5.2.27}
    \label{eq:5.2.27}
\end{equation}

\eqref{eq:5.2.25} 与 \eqref{eq:5.2.27} 是同一个一阶线性方程，并具有相同的无穷远归一化，因此由唯一性

\begin{equation}
    \mu_j^H(x,t,\bar z)
    =
    \mu_j^{-1}(x,t,z),
    \qquad j=1,2.
    \tag{5.2.28}
    \label{eq:5.2.28}
\end{equation}

将散射关系改写为

\begin{equation}
    S(z)
    =
    e^{i\theta(z)\widehat{\sigma}_3}
    \left[\mu_2^{-1}(x,t,z)\mu_1(x,t,z)\right],
    \tag{5.2.29}
    \label{eq:5.2.29}
\end{equation}

再利用 \eqref{eq:5.2.28}，得到

\[
    S^H(\bar z)
    =
    \begin{pmatrix}
        \overline{s_{11}(\bar z)} & \overline{s_{21}(\bar z)} \\
        \overline{s_{12}(\bar z)} & \overline{s_{22}(\bar z)}
    \end{pmatrix}
    =
    \begin{pmatrix}
        s_{22}(z)  & -s_{12}(z) \\
        -s_{21}(z) & s_{11}(z)
    \end{pmatrix}
    =
    S^{-1}(z),
\]

从而矩阵元满足相应的共轭关系，可写为

\begin{equation}
    \overline{s_{11}(\bar z)}=s_{22}(z),
    \qquad
    \overline{s_{12}(\bar z)}=-s_{21}(z).
    \tag{5.2.30}
    \label{eq:5.2.30}
\end{equation}

\begin{callout}[2.2 节真正要记住什么？]
    不是每一步共轭转置计算，而是以下逻辑：

    \[
        \begin{aligned}
             & P^H=-P
            \Longrightarrow
            \mu_j^H(\bar z)\text{ 与 }\mu_j^{-1}(z)
            \text{ 满足同一方程和同一归一化},
            \\
             & \mu_j^H(\bar z)=\mu_j^{-1}(z)
            \Longrightarrow
            S^H(\bar z)=S^{-1}(z).
        \end{aligned}
    \]
    这个对称性后面保证上半平面的离散零点与下半平面的共轭零点成对出现。
\end{callout}

\section{构造并求解 RH 问题}

\subsection{规范化 RH 问题}

基于第 5.2 节得到的解析列和解析行，引入投影矩阵

\begin{equation}
    H_1=
    \begin{pmatrix}
        1 & 0 \\
        0 & 0
    \end{pmatrix},
    \qquad
    H_2=
    \begin{pmatrix}
        0 & 0 \\
        0 & 1
    \end{pmatrix}.
    \tag{5.3.1}
    \label{eq:5.3.1}
\end{equation}

右乘 $H_1,H_2$ 分别选取第一、二列；左乘 $H_1,H_2$ 分别选取第一、二行。因此定义

\[
    P_+
    =
    \mu_1H_1+\mu_2H_2
    =
    (\mu_1^+,\mu_2^+),
\]

\[
    P_-
    =
    H_1\mu_1^{-1}+H_2\mu_2^{-1}
    =
    \begin{pmatrix}
        \widehat\mu_1^- \\
        \widehat\mu_2^-
    \end{pmatrix}.
\]

由第 5.2 节的解析性，$P_+$ 在 $\mathbb{C}^+$ 解析，$P_-$ 在 $\mathbb{C}^-$ 解析，而且

\[
    P_\pm(x,t,z)\to I,\qquad z\to\infty.
\]

由 \eqref{eq:5.2.28} 还得到对称性

\begin{equation}
    P_+^H(x,t,\bar z)
    =
    H_1\mu_1^H(x,t,\bar z)
    +
    H_2\mu_2^H(x,t,\bar z)
    =
    P_-(x,t,z).
    \tag{5.3.3}
    \label{eq:5.3.3}
\end{equation}

于是得到本章真正需要的 RH 问题：

\begin{equation}
    P_\pm(x,t,z)\text{ 在 }\mathbb{C}^\pm\text{ 上解析},
    \tag{5.3.4}
    \label{eq:5.3.4}
\end{equation}

\begin{equation}
    P_-(x,t,z)P_+(x,t,z)
    =
    G(x,t,z),
    \qquad z\in\mathbb{R},
    \tag{5.3.5}
    \label{eq:5.3.5}
\end{equation}

\begin{equation}
    P_\pm(x,t,z)\to I,
    \qquad z\to\infty.
    \tag{5.3.6}
    \label{eq:5.3.6}
\end{equation}

跳跃矩阵为

\begin{equation}
    G(x,t,z)
    =
    e^{-i\theta(z)\widehat{\sigma}_3}
    \begin{pmatrix}
        1         & -s_{12}(z) \\
        s_{21}(z) & 1
    \end{pmatrix}=
    \begin{pmatrix}
        1                     & -s_{12}(z)e^{-2i\theta} \\
        s_{21}(z)e^{2i\theta} & 1
    \end{pmatrix}.
    \tag{5.3.7}
    \label{eq:5.3.7}
\end{equation}

NLS 的解可由 RH 解重构：

\begin{equation}
    q(x,t)
    =
    2i\lim\limits_{z\to\infty}(zP_+)_{12}
    =
    2i(P_+^{(1)})_{12},
    \tag{5.3.8}
    \label{eq:5.3.8}
\end{equation}

其中

\begin{equation}
    P_+(x,t,z)
    =
    I+\dfrac{P_+^{(1)}(x,t)}{z}
    +O(z^{-2}).
    \tag{5.3.9}
    \label{eq:5.3.9}
\end{equation}

\begin{callout}[计算：跳跃矩阵 $G$ 从哪里来？]
    由定义

    \[
        P_-P_+
        =
        (H_1\mu_1^{-1}+H_2\mu_2^{-1})
        (\mu_1H_1+\mu_2H_2).
    \]
    完全展开：

    \[
        \begin{aligned}
            P_-P_+
            ={} &
            H_1\mu_1^{-1}\mu_1H_1
            +
            H_1\mu_1^{-1}\mu_2H_2 \\
                & +
            H_2\mu_2^{-1}\mu_1H_1
            +
            H_2\mu_2^{-1}\mu_2H_2.
        \end{aligned}
    \]
    第一、四项给出

    \[
        H_1+H_2=I.
    \]
    由散射关系 \eqref{eq:5.2.3}，

    \[
        \mu_2^{-1}\mu_1
        =
        e^{-i\theta\widehat{\sigma}_3}S,
    \]

    \[
        \mu_1^{-1}\mu_2
        =
        e^{-i\theta\widehat{\sigma}_3}S^{-1}.
    \]
    又因为 $\det S=1$，

    \[
        S^{-1}
        =
        \begin{pmatrix}
            s_{22}  & -s_{12} \\
            -s_{21} & s_{11}
        \end{pmatrix}.
    \]
    利用 $H_1,H_2$ 只保留相应行列，最后得到

    \[
        P_-P_+
        =
        e^{-i\theta\widehat{\sigma}_3}
        \begin{pmatrix}
            1      & -s_{12} \\
            s_{21} & 1
        \end{pmatrix},
    \]
    即 \eqref{eq:5.3.7}。
    因为

    \[
        e^{-i\theta\widehat{\sigma}_3}
        \begin{pmatrix}
            a & b \\c&d
        \end{pmatrix}
        =
        \begin{pmatrix}
            a             & e^{-2i\theta}b \\
            e^{2i\theta}c & d
        \end{pmatrix},
    \]
    所以也可以把跳跃矩阵理解为

    \[
        G=
        \begin{pmatrix}
            1                     & -s_{12}(z)e^{-2i\theta} \\
            s_{21}(z)e^{2i\theta} & 1
        \end{pmatrix}.
    \]
    此外，直接由散射关系可得两个非常重要的结果：

    \[
        \det P_+(z)=s_{11}(z),
        \qquad
        \det P_-(z)=s_{22}(z).
    \]
    以 $P_+$ 为例，由

    \[
        \mu_1
        =
        \mu_2e^{-i\theta\widehat{\sigma}_3}S
    \]
    可写成

    \[
        P_+
        =
        \mu_1H_1+\mu_2H_2
        =
        \mu_2
        \left[
            \left(e^{-i\theta\widehat{\sigma}_3}S\right)H_1+H_2
            \right].
    \]
    而

    \[
        e^{-i\theta\widehat{\sigma}_3}S
        =
        \begin{pmatrix}
            s_{11}             & s_{12}e^{-2i\theta} \\
            s_{21}e^{2i\theta} & s_{22}
        \end{pmatrix},
    \]
    所以

    \[
        \left(e^{-i\theta\widehat{\sigma}_3}S\right)H_1+H_2
        =
        \begin{pmatrix}
            s_{11}             & 0 \\
            s_{21}e^{2i\theta} & 1
        \end{pmatrix}.
    \]
    因此

    \[
        \det P_+
        =
        \det\mu_2\cdot s_{11}
        =
        s_{11}.
    \]
    同理可得 $\det P_-=s_{22}$。
\end{callout}

\begin{callout}[到这里发生了什么？]
    第 5.1 节把 $q$ 与特征函数的大 $z$ 系数联系起来（\eqref{eq:5.1.15}）：

    \[
        q(x,t)=2i(\mu^{(1)})_{12}=2i\lim\limits_{z\to\infty}(z\mu)_{12}.
    \]
    第 5.2 节得到特征函数的解析性与散射数据；第 5.3.1 节把这些解析列、解析行重新拼装成 $P_\pm$，于是终于得到

    \[
        \boxed{\text{NLS 初值问题}\longrightarrow\text{RH 问题}.}
    \]
    恢复 $q$ 的公式也随之从 $\mu$ 的版本换成 $P_+$ 的版本（\eqref{eq:5.3.8}）：

    \[
        q(x,t)=2i\lim\limits_{z\to\infty}(zP_+)_{12}=2i(P_+^{(1)})_{12}.
    \]
    后面只需要解决这个 RH 问题，再代回 \eqref{eq:5.3.8} 恢复 $q$。
\end{callout}

\subsection{RH 问题的可解性}

\subsubsection{正则 RH 问题}

如果

\begin{equation}
    \det P_\pm(z)\ne0,
    \tag{5.3.10}
    \label{eq:5.3.10}
\end{equation}

即相应解析区域中不存在零点，则称 RH 问题为正则的。由 \eqref{eq:5.3.5}

\[
    P_-P_+=G
\]

右乘 $P_+^{-1}$ 后得到

\[
    P_-=GP_+^{-1}.
\]

因此

\[
    P_+^{-1}-P_-
    =
    (I-G)P_+^{-1}
    \triangleq
    \widehat G\,P_+^{-1}.
\]

这是一个\textbf{加性 RH 跳跃关系}。由 Plemelj(普莱梅尔) 公式，

\[
    P_+^{-1}(z)
    =
    I+
    \dfrac{1}{2\pi i}
    \int_{-\infty}^{\infty}
    \dfrac{\widehat G(s)P_+^{-1}(s)}{s-z}\,\mathrm ds,
    \qquad
    z\in\mathbb{C}^+.
\]

注意右端仍含 $P_+^{-1}$，所以这不是显式解，而是把 RH 问题转化成了一个奇异积分方程。

\begin{callout}[前置工具：这里具体用了哪个 Plemelj 公式？（原书 3.5.3 节，(3.5.26)---(3.5.28)）]
    前面复分析部分的加性 RH 问题是

    \[
        F_+-F_-=f,
    \]
    其 Cauchy 表示为

    \[
        F(z)
        =
        A+
        \dfrac{1}{2\pi i}
        \int\limits_\Sigma
        \dfrac{f(\xi)}{\xi-z}\,\mathrm d\xi.
    \]
    本节取

    \[
        F_+=P_+^{-1},
        \qquad
        F_-=P_-,
        \qquad
        f=\widehat G\,P_+^{-1},
    \]
    跳跃曲线取 $\Sigma=\mathbb{R}$，再由 $P_+^{-1}\to I$ 取 $A=I$，就得到上面的积分方程。
    所以这里的逻辑是：

    \[
        \text{乘法跳跃}
        \longrightarrow
        \text{加法跳跃}
        \longrightarrow
        \text{Plemelj}
        \longrightarrow
        \text{奇异积分方程}.
    \]

\end{callout}

\subsubsection{非正则 RH 问题}

若 \eqref{eq:5.3.10} 不成立，假设 $z_j\in\mathbb{C}^+$ 是 $\det P_+=s_{11}$ 的离散零点，则由对称性，$\bar z_j\in\mathbb{C}^-$ 是 $\det P_-=s_{22}$ 的对应零点。
因为矩阵在这些点奇异，所以存在非零列向量、行向量

\begin{equation}
    P_+(z_j)w_j=0,
    \qquad j=1,\ldots,N,
    \tag{5.3.11}
    \label{eq:5.3.11}
\end{equation}

\begin{equation}
    w_j^*P_-(\bar z_j)=0.
    \tag{5.3.12}
    \label{eq:5.3.12}
\end{equation}

由 \eqref{eq:5.3.3} 可取

\begin{equation}
    w_j^*=w_j^H.
    \tag{5.3.13}
    \label{eq:5.3.13}
\end{equation}

\begin{callout}[概念：为什么零点对应``离散谱''？（原书 5.3.2 节）]
    $\det P_+(z)=s_{11}(z)$。
    当 $s_{11}(z_j)=0$ 时，$P_+(z_j)$ 失去可逆性，出现非平凡核向量 $w_j$。
    这些不位于连续谱实轴上的孤立复零点就是离散谱。
    后面孤子正是由这些离散谱数据产生。
\end{callout}

\textbf{定理 5.3（Zakharov--Shabat 零点剥离）}\\
带有上述简单零点结构的非正则 RH 问题可分解为

\begin{equation}
    P_+(z)
    =
    \widehat P_+(z)\Gamma(z),
    \tag{5.3.14}
    \label{eq:5.3.14}
\end{equation}

\begin{equation}
    P_-(z)
    =
    \Gamma^{-1}(z)\widehat P_-(z),
    \tag{5.3.15}
    \label{eq:5.3.15}
\end{equation}

其中

\begin{equation}
    \Gamma(z)
    =
    I+
    \sum\limits_{k,j=1}^{N}
    \dfrac{
        w_k(M^{-1})_{kj}w_j^*
    }{
        z-\bar z_j
    },
    \tag{5.3.16}
    \label{eq:5.3.16}
\end{equation}

\begin{equation}
    \Gamma^{-1}(z)
    =
    I-
    \sum\limits_{k,j=1}^{N}
    \dfrac{
        w_k(M^{-1})_{kj}w_j^*
    }{
        z-z_k
    },
    \tag{5.3.17}
    \label{eq:5.3.17}
\end{equation}

而

\[
    M_{kj}
    =
    \dfrac{w_k^*w_j}{\bar z_k-z_j},
    \qquad
    1\le k,j\le N.
\]

并且

\[
    \det\Gamma(z)
    =
    \prod\limits_{j=1}^{N}
    \dfrac{z-z_j}{z-\bar z_j}.
\]

剥离后的 $\widehat P_\pm$ 满足正则 RH 问题：

\begin{itemize}
    \item $\widehat P_\pm$ 在 $\mathbb{C}^\pm$ 上解析；
    \item 在实轴上

          \[
              \widehat P_-(z)\widehat P_+(z)
              =
              \Gamma(z)G(z)\Gamma^{-1}(z);
          \]

    \item $\widehat P_\pm(z)\to I$，$z\to\infty$。
\end{itemize}

核心思想就是：\textbf{把所有离散零点结构装入有理矩阵 $\Gamma$，剩余的 $\widehat P_\pm$ 不再含这些零点，于是非正则 RH 问题被化归为正则 RH 问题。}

\begin{callout}[定理 5.3 的证明思路（组会可略）]
    对第 $j$ 个零点引入单极点因子

    \[
        \Gamma_j(z)
        =
        I+
        \dfrac{\bar z_j-z_j}{z-\bar z_j}
        \dfrac{v_jv_j^*}{v_j^*v_j},
    \]

    \begin{equation}
        \Gamma_j^{-1}(z)
        =
        I-
        \dfrac{\bar z_j-z_j}{z-z_j}
        \dfrac{v_jv_j^*}{v_j^*v_j}.
        \tag{5.3.19}
        \label{eq:5.3.19}
    \end{equation}
    逐个消去 $P_+$ 在 $z_j$ 和 $P_-$ 在 $\bar z_j$ 的零点。核向量与逐步变换后的向量满足

    \begin{equation}
        w_j(z_j)
        =
        \Gamma_1^{-1}(z_j)\cdots
        \Gamma_{j-1}^{-1}(z_j)v_j(z_j),
        \tag{5.3.18}
        \label{eq:5.3.18}
    \end{equation}
    以及相应的左核关系。最后令

    \begin{equation}
        \Gamma(z)
        =
        \Gamma_N(z)\cdots\Gamma_1(z),
        \qquad
        \Gamma^{-1}(z)
        =
        \Gamma_1^{-1}(z)\cdots\Gamma_N^{-1}(z).
        \tag{5.3.20}
        \label{eq:5.3.20}
    \end{equation}
    为得到 \eqref{eq:5.3.16}---\eqref{eq:5.3.17} 的显式形式，再作部分分式展开

    \begin{equation}
        \Gamma(z)
        =
        I+
        \sum\limits_{j=1}^{N}
        \dfrac{\zeta_jw_j^*}{z-\bar z_j},
        \tag{5.3.21}
        \label{eq:5.3.21}
    \end{equation}

    \begin{equation}
        \Gamma^{-1}(z)
        =
        I-
        \sum\limits_{j=1}^{N}
        \dfrac{w_j\zeta_j^*}{z-z_j}.
        \tag{5.3.22}
        \label{eq:5.3.22}
    \end{equation}
    由 $\Gamma\Gamma^{-1}=I$，在 $z=z_k$ 取留数：

    \begin{equation}
        0
        =
        -\Gamma(z_k)w_k\zeta_k^*
        =
        \left(
        -w_k+
        \sum\limits_{j=1}^{N}
        \dfrac{w_j^*w_k}{\bar z_j-z_k}\zeta_j
        \right)\zeta_k^*.
        \tag{5.3.23}
        \label{eq:5.3.23}
    \end{equation}
    因而

    \begin{equation}
        -w_k+
        \sum\limits_{j=1}^{N}
        \dfrac{w_j^*w_k}{\bar z_j-z_k}\zeta_j
        =
        0.
        \tag{5.3.24}
        \label{eq:5.3.24}
    \end{equation}
    写成分块线性方程组：

    \begin{equation}
        (\zeta_1,\ldots,\zeta_N)M
        =
        (w_1,\ldots,w_N).
        \tag{5.3.25}
        \label{eq:5.3.25}
    \end{equation}
    解得 $\displaystyle \zeta_j=\sum\limits_k w_k(M^{-1})_{kj}$，代回即得到 \eqref{eq:5.3.16}---\eqref{eq:5.3.17}。
\end{callout}

\section{NLS 方程的 \texorpdfstring{$N$}{N} 孤子解}

\subsection{离散核向量的时空演化}

从 \eqref{eq:5.3.11}

\[
    P_+(z_j)w_j=0
\]

分别对 $x,t$ 求导：

\begin{equation}
    P_{+,x}w_j+P_+w_{j,x}=0,
    \tag{5.4.1}
    \label{eq:5.4.1}
\end{equation}

\begin{equation}
    P_{+,t}w_j+P_+w_{j,t}=0.
    \tag{5.4.2}
    \label{eq:5.4.2}
\end{equation}

利用 $P_+$ 的定义和 Lax 对 \eqref{eq:5.1.7}，得到

\begin{equation}
    P_{+,x}
    =
    \mu_{1,x}H_1+\mu_{2,x}H_2
    =
    -iz_j[\sigma_3,P_+]+PP_+.
    \tag{5.4.3}
    \label{eq:5.4.3}
\end{equation}

同理由 \eqref{eq:5.1.8} 得到

\begin{equation}
    P_{+,t}
    =
    -2iz_j^2[\sigma_3,P_+]+QP_+.
    \tag{5.4.4}
    \label{eq:5.4.4}
\end{equation}

代入 \eqref{eq:5.4.1}---\eqref{eq:5.4.2}，再使用 $P_+(z_j)w_j=0$：

\[
    P_+\left(w_{j,x}+iz_j\sigma_3w_j\right)=0,
\]

\[
    P_+\left(w_{j,t}+2iz_j^2\sigma_3w_j\right)=0.
\]

选择核向量的规范，使括号内向量为零：

\[
    w_{j,x}+iz_j\sigma_3w_j=0,
    \qquad
    w_{j,t}+2iz_j^2\sigma_3w_j=0.
\]

解得

\begin{equation}
    w_j
    =
    e^{-i\theta(z_j)\sigma_3}w_{j0},
    \qquad
    j=1,\ldots,N,
    \tag{5.4.5}
    \label{eq:5.4.5}
\end{equation}

其中 $w_{j0}$ 为二维常数列向量，因此

\begin{equation}
    w_j^*
    =
    w_{j0}^He^{i\theta(\bar z_j)\sigma_3}.
    \tag{5.4.6}
    \label{eq:5.4.6}
\end{equation}

\subsection{\texorpdfstring{$N$}{N} 孤子解公式}

先从剥离零点后的正则 RH 问题出发：

\[
    \widehat P_-(z)\widehat P_+(z)
    =
    \Gamma(z)G(z)\Gamma^{-1}(z),
    \qquad z\in\mathbb R.
\]

右乘 $\widehat P_+^{-1}$，再移项，得到加性跳跃关系

\[
    \begin{aligned}
        \widehat P_+^{-1}-\widehat P_-
        & =
        \left[I-\Gamma G\Gamma^{-1}\right]\widehat P_+^{-1} \\
        & =
        \Gamma(I-G)\Gamma^{-1}\widehat P_+^{-1} \\
        & =
        \Gamma\widehat G\Gamma^{-1}\widehat P_+^{-1}.
    \end{aligned}
\]

因此由 Plemelj 公式，

\[
    \widehat P_+^{-1}(z)
    =
    I+
    \dfrac{1}{2\pi i}
    \int_{-\infty}^{\infty}
    \dfrac{
        \Gamma(s)\widehat G(s)\Gamma^{-1}(s)\widehat P_+^{-1}(s)
    }{s-z}
    \,\mathrm ds.
\]

当 $z\to\infty$ 时，

\[
    \dfrac{1}{s-z}
    =
    -\dfrac{1}{z}\dfrac{1}{1-s/z}
    =
    -\dfrac{1}{z}
    \left(1+\dfrac{s}{z}+O(z^{-2})\right),
\]

所以

\[
    \widehat P_+^{-1}(z)
    =
    I-
    \dfrac{1}{2\pi iz}
    \int_{-\infty}^{\infty}
    \Gamma\widehat G\Gamma^{-1}\widehat P_+^{-1}
    \,\mathrm ds
    +O(z^{-2}).
\]

若记

\[
    A=
    \dfrac{1}{2\pi i}
    \int_{-\infty}^{\infty}
    \Gamma\widehat G\Gamma^{-1}\widehat P_+^{-1}
    \,\mathrm ds,
\]

则

\[
    \widehat P_+^{-1}(z)=I-\dfrac{A}{z}+O(z^{-2}).
\]

利用

\[
    \left(I+\dfrac{A}{z}+O(z^{-2})\right)
    \left(I-\dfrac{A}{z}+O(z^{-2})\right)
    =I+O(z^{-2}),
\]

可得

\begin{equation}
    \widehat P_+(z)
    =
    I+
    \dfrac{1}{2\pi iz}
    \int_{-\infty}^{\infty}
    \Gamma\widehat G\Gamma^{-1}\widehat P_+^{-1}\,\mathrm ds
    +
    O(z^{-2}).
    \tag{5.4.7}
    \label{eq:5.4.7}
\end{equation}

另一方面，由 \eqref{eq:5.3.16}

\begin{equation}
    \Gamma(z)
    =
    I+
    \dfrac{1}{z}
    \sum\limits_{k,j=1}^{N}
    w_k(M^{-1})_{kj}w_j^*
    +
    O(z^{-2}).
    \tag{5.4.8}
    \label{eq:5.4.8}
\end{equation}

利用

\[
    P_+=\widehat P_+\Gamma
\]

以及 \eqref{eq:5.3.9}，比较 $z^{-1}$ 系数：

\begin{equation}
    P_+^{(1)}
    =
    \dfrac{1}{2\pi i}
    \int_{-\infty}^{\infty}
    \Gamma\widehat G\Gamma^{-1}\widehat P_+^{-1}\,\mathrm ds
    +
    \sum\limits_{k,j=1}^{N}
    w_k(M^{-1})_{kj}w_j^*.
    \tag{5.4.9}
    \label{eq:5.4.9}
\end{equation}

这个式子非常重要：第一项是连续谱贡献，第二项是离散谱贡献。
特别地，当散射数据

\[
    s_{12}=s_{21}=0
\]

时，由 \eqref{eq:5.3.7} 有 $G=I,\widehat G=I-G=0$，连续谱积分项消失：

\[
    P_+^{(1)}
    =
    \sum\limits_{k,j=1}^{N}
    w_k(M^{-1})_{kj}w_j^*.
\]

记

\[
    \lambda_j
    =
    -i(z_jx+2z_j^2t),
\]

并取

\[
    w_{j0}
    =
    \begin{pmatrix}
        c_j \\1
    \end{pmatrix}.
\]

由 \eqref{eq:5.4.5}

\begin{equation}
    w_j
    =
    \begin{pmatrix}
        e^{\lambda_j} & 0              \\
        0             & e^{-\lambda_j}
    \end{pmatrix}
    \begin{pmatrix}
        c_j \\1
    \end{pmatrix}
    =
    \begin{pmatrix}
        c_je^{\lambda_j} \\
        e^{-\lambda_j}
    \end{pmatrix}.
    \tag{5.4.10}
    \label{eq:5.4.10}
\end{equation}

因此

\begin{equation}
    w_j^*
    =
    \left(
    \bar c_je^{\bar\lambda_j},
    e^{-\bar\lambda_j}
    \right).
    \tag{5.4.11}
    \label{eq:5.4.11}
\end{equation}

矩阵 $M$ 的元素为

\begin{equation}
    M_{kj}
    =
    \dfrac{w_k^*w_j}{\bar z_k-z_j}
    =
    \dfrac{
        \bar c_kc_je^{\bar\lambda_k+\lambda_j}
        +
        e^{-\bar\lambda_k-\lambda_j}
    }{
        \bar z_k-z_j
    }.
    \tag{5.4.12}
    \label{eq:5.4.12}
\end{equation}

由重构公式 \eqref{eq:5.3.8}：

\begin{equation}
    \begin{aligned}
        q
        & =
        2i(P_+^{(1)})_{12} \\
        & =
        2i\left(
        \sum\limits_{k,j=1}^{N}
        w_k(M^{-1})_{kj}w_j^*
        \right)_{12} \\
        & =
        2i\left(
        \sum\limits_{k,j=1}^{N}
        (M^{-1})_{kj}
        \begin{pmatrix}
            c_k\bar c_j e^{\lambda_k+\bar\lambda_j}
            & c_ke^{\lambda_k-\bar\lambda_j} \\
            \bar c_j e^{-\lambda_k+\bar\lambda_j}
            & e^{-\lambda_k-\bar\lambda_j}
        \end{pmatrix}
        \right)_{12} \\
        & =
        2i
        \sum\limits_{k,j=1}^{N}
        c_ke^{\lambda_k-\bar\lambda_j}(M^{-1})_{kj}.
    \end{aligned}
    \tag{5.4.13}
    \label{eq:5.4.13}
\end{equation}

定义 $(N+1)\times(N+1)$ 矩阵

\begin{equation}
    R=
    \begin{pmatrix}
        0                  & c_1e^{\lambda_1} & \cdots & c_Ne^{\lambda_N} \\
        e^{-\bar\lambda_1} & M_{11}           & \cdots & M_{1N}           \\
        \vdots             & \vdots           &        & \vdots           \\
        e^{-\bar\lambda_N} & M_{N1}           & \cdots & M_{NN}
    \end{pmatrix}.
    \tag{5.4.14}
    \label{eq:5.4.14}
\end{equation}

按第一行、第一列展开，或者等价地使用 Schur 补，可把双重求和写成伴随矩阵形式：

\begin{equation}
    \det R
    =
    -\sum\limits_{k,j=1}^{N}
    c_ke^{\lambda_k-\bar\lambda_j}(M^*)_{jk}
    =
    -\det M
    \sum\limits_{k,j=1}^{N}
    c_ke^{\lambda_k-\bar\lambda_j}(M^{-1})_{kj}.
    \tag{5.4.15}
    \label{eq:5.4.15}
\end{equation}

这里 $M^*$ 在这一式中表示 \textbf{伴随矩阵（adjugate）}。使用

\[
    M^{-1}
    =
    \dfrac{1}{\det M}M^*,
    \qquad
    (M^{-1})_{kj}
    =
    \dfrac{(M^*)_{jk}}{\det M}.
\]

于是

\begin{equation}
    \sum\limits_{k,j=1}^{N}
    c_ke^{\lambda_k-\bar\lambda_j}(M^{-1})_{kj}
    =
    -\dfrac{\det R}{\det M}.
    \tag{5.4.16}
    \label{eq:5.4.16}
\end{equation}

代回 \eqref{eq:5.4.13}，得到 $N$ 孤子解

\begin{equation}
    \boxed{
        q(x,t)
        =
        -2i\dfrac{\det R}{\det M}.
    }
    \tag{5.4.17}
    \label{eq:5.4.17}
\end{equation}

\begin{callout}[4.2 节的核心思路]
    一般 RH 解的 $z^{-1}$ 系数由两部分组成：

    \[
        P_+^{(1)}
        =
        \text{连续谱积分}
        +
        \text{离散谱有限和}.
    \]
    无反射 $s_{12}=s_{21}=0$ 时，$G=I$，连续谱积分消失，只剩离散谱有限和；再借助有限维矩阵 $M$ 和 $R$，就得到行列式形式 \eqref{eq:5.4.17}。因此：

    \[
        \boxed{\text{连续谱对应辐射/反射，离散谱对应孤子}.}
    \]

\end{callout}

\subsection{单孤子解}

现在令 $N=1$。取一对离散谱点

\[
    z_1\in\mathbb{C}^+,
    \qquad
    \bar z_1\in\mathbb{C}^-,
\]

并记

\[
    \lambda_1
    =
    -i(z_1x+2z_1^2t).
\]

由 \eqref{eq:5.4.10}---\eqref{eq:5.4.12}，

\[
    w_1
    =
    \begin{pmatrix}
        c_1e^{\lambda_1} \\
        e^{-\lambda_1}
    \end{pmatrix},
    \qquad
    w_1^*
    =
    \left(
    \bar c_1e^{\bar\lambda_1},
    e^{-\bar\lambda_1}
    \right).
\]

此时 $M$ 是 $1\times1$ 矩阵，所以

\[
    \det M
    =
    M_{11}
    =
    \dfrac{
        \lvert c_1\rvert^2e^{\lambda_1+\bar\lambda_1}
        +
        e^{-(\lambda_1+\bar\lambda_1)}
    }{
        \bar z_1-z_1
    }.
\]

而

\[
    R=
    \begin{pmatrix}
        0                  & c_1e^{\lambda_1} \\
        e^{-\bar\lambda_1} & M_{11}
    \end{pmatrix},
\]

因此

\[
    \det R
    =
    -c_1e^{\lambda_1-\bar\lambda_1}.
\]

代入 \eqref{eq:5.4.17}，得到单孤子解

\begin{equation}
    q(x,t)
    =
    2i(\bar z_1-z_1)
    \dfrac{
        c_1e^{\lambda_1-\bar\lambda_1}
    }{
        \lvert c_1\rvert^2e^{\lambda_1+\bar\lambda_1}
        +
        e^{-(\lambda_1+\bar\lambda_1)}
    }.
    \tag{5.4.18}
    \label{eq:5.4.18}
\end{equation}

进一步取

\[
    z_1=\xi+i\eta,
    \qquad
    c_1=e^{-\delta_0+ik_0},
    \qquad
    \eta>0.
\]

由

\[
    \lambda_1=-i(z_1x+2z_1^2t)
\]

可得

\[
    \lambda_1+\bar\lambda_1
    =
    2\eta(x+4\xi t),
\]

\[
    \lambda_1-\bar\lambda_1
    =
    -2i\xi x-4i(\xi^2-\eta^2)t,
\]

以及

\[
    2i(\bar z_1-z_1)=4\eta.
\]

代入 \eqref{eq:5.4.18} 并化简：

\[
    \boxed{
        q(x,t)
        =
        2\eta\,
        e^{-2i\xi x-4i(\xi^2-\eta^2)t+ik_0}
        \operatorname{sech}
        \!\left[
            2\eta(x+4\xi t)-\delta_0
            \right].
    }
\]

这个解的模长是局域的 $\operatorname{sech}$ 型波包，因此属于零边界下的亮孤子。

\begin{callout}[参数的直观意义]
    $z_1=\xi+i\eta$ 中，$\eta>0$ 控制孤子的振幅和宽度，$\xi$ 控制传播速度和相位；$c_1=e^{-\delta_0+ik_0}$ 中，$\delta_0$ 与孤子的位置平移有关，$k_0$ 是额外相位。
    离散谱点决定孤子的主要动力学参数，而归一化常数决定位置、相位等自由度。
\end{callout}

\section*{最后的逻辑闭环}

现在可以把整章压缩成四个层次。

\begin{enumerate}
    \item \textbf{第一层：非线性方程与线性谱问题}

          \[
              \text{NLS}
              \overset{\text{Lax 对}}{\longleftrightarrow}
              \text{线性谱问题}.
          \]

    \item \textbf{第二层：谱问题与散射数据}

          \[
              \text{Jost 解}
              \longrightarrow
              S(z)
              \longrightarrow
              \begin{cases}
                  \text{连续谱数据 }s_{12},s_{21}, \\
                  \text{离散谱数据 }z_j,w_j.
              \end{cases}
          \]

    \item \textbf{第三层：谱数据与 RH 问题}

          \[
              \text{解析列/解析行}
              +
              \text{散射数据}
              \longrightarrow
              (P_+,P_-,G).
          \]

          正则时用 Plemelj；非正则时先用 $\Gamma$ 剥离零点。

    \item \textbf{第四层：RH 解回到 NLS}

          \[
              P_+(z)
              =
              I+\dfrac{P_+^{(1)}}{z}+O(z^{-2})
              \longrightarrow
              q(x,t)=2i(P_+^{(1)})_{12}.
          \]

          无反射时 $G=I$，只剩离散谱：

          \[
              q(x,t)
              =
              -2i\dfrac{\det R}{\det M}.
          \]
\end{enumerate}

\begin{callout}[组会最后一句]
    这一章真正建立的是一条完整的逆散射/RH 链条：

    \[
        \boxed{
            q_0
            \longrightarrow
            \text{Jost 解}
            \longrightarrow
            \text{散射数据}
            \longrightarrow
            \text{RH 问题}
            \longrightarrow
            P_+
            \longrightarrow
            q(x,t).
        }
    \]
    Lax 对负责把非线性 PDE 编码成谱问题；Jost 解和散射矩阵负责提取谱数据；解析性负责组织 RH 问题；Plemelj 公式和 $\Gamma$ 零点剥离负责求解 RH 问题；最后用大谱参数展开完成重构。
    无反射条件去掉连续谱贡献后，离散谱直接产生孤子。
\end{callout}


\end{document}
